% ===== PRÉAMBULE CANONIQUE — à coller VERBATIM en tête de CHAQUE .tex =====
\documentclass[11pt,a4paper]{article}
\usepackage[utf8]{inputenc}\usepackage[T1]{fontenc}\usepackage{lmodern}
\usepackage[french]{babel}
\usepackage{amsmath,amssymb,mathtools}\usepackage{icomma}
\usepackage[version=4]{mhchem}          % \ce{...} : équations & formules
\usepackage{siunitx}\sisetup{locale=FR} % \qty{}{}, \si{}, \ang{}
\usepackage{chemfig}                    % \chemfig{...} : structures développées
\usepackage{xcolor}\usepackage{colortbl}
\usepackage{tikz}\usepackage{pgfplots}\pgfplotsset{compat=1.18}
\usepackage[most]{tcolorbox}\usepackage{enumitem}\usepackage{array,booktabs}
\usepackage[a4paper,margin=2.2cm]{geometry}
\usetikzlibrary{arrows.meta,calc,angles,quotes,positioning,patterns,backgrounds,shapes.geometric}
\definecolor{primary}{RGB}{222,49,99}\definecolor{primarydark}{RGB}{150,22,55}
\definecolor{defcol}{RGB}{0,114,178}\definecolor{methcol}{RGB}{52,92,148}
\definecolor{excol}{RGB}{0,135,90}\definecolor{attncol}{RGB}{200,40,55}
\tcbset{boxbase/.style={enhanced,breakable,boxrule=0.9pt,arc=2.5pt,
  left=3mm,right=3mm,top=2.3mm,bottom=2.3mm,fonttitle=\bfseries,coltitle=white}}
% --- boîtes TUTORIEL (noms EXACTS, ne pas renommer) ---
\newtcolorbox{cle}[1][]{boxbase,colback=defcol!6,colframe=defcol,
  title={\textsc{À retenir}\ifx\relax#1\relax\relax\else\textmd{ : \itshape #1}\fi}}
\newtcolorbox{methode}[1][]{boxbase,colback=methcol!7,colframe=methcol,sharp corners,
  title={\textsc{Méthode}\ifx\relax#1\relax\relax\else\textmd{ --- \itshape #1}\fi}}
\newtcolorbox{exemplebox}[1][]{boxbase,colback=excol!5,colframe=excol,
  title={\textsc{Exemple}\ifx\relax#1\relax\relax\else\textmd{ : \itshape #1}\fi}}
\newtcolorbox{piege}[1][]{boxbase,colback=attncol!6,colframe=attncol,
  title={\textsc{Piège}\ifx\relax#1\relax\relax\else\textmd{ : \itshape #1}\fi}}
\newtcolorbox{remarque}[1][]{boxbase,colback=black!4,colframe=black!45,
  title={\textsc{Remarque}\ifx\relax#1\relax\relax\else\textmd{ : \itshape #1}\fi}}
% --- boîtes EXERCICE / CORRIGÉ (noms EXACTS, auto-comptées) ---
\newtcolorbox[auto counter]{exo}[1][]{boxbase,colback=primary!4,colframe=primary,
  title={\textsc{Exercice}~\thetcbcounter\ifx\relax#1\relax\relax\else\textmd{ --- \itshape #1}\fi}}
\newtcolorbox[auto counter]{corrige}[1][]{boxbase,colback=excol!4,colframe=excol,
  title={\textsc{Corrigé}~\thetcbcounter\ifx\relax#1\relax\relax\else\textmd{ --- \itshape #1}\fi}}
\newcommand{\R}{\mathbb{R}}\newcommand{\N}{\mathbb{N}}\newcommand{\Z}{\mathbb{Z}}\newcommand{\ds}{\displaystyle}
% (un \usepackage{fancyhdr}+en-tête est possible ; il est ignoré côté web)
% =========================================================================
\begin{document}
\begin{corrige}[classer par électronégativité]
On combine les deux sens : vers la droite \emph{et} vers le haut. $\ce{F}$ est au coin haut-droite
(le plus électronégatif) ; $\ce{N}$ est sur la même période mais plus à gauche ; $\ce{Al}$ est plus bas
et à gauche ; $\ce{K}$ est en bas à gauche (le moins électronégatif). D'où :
\[ \ce{F} > \ce{N} > \ce{Al} > \ce{K}. \]
\end{corrige}

\begin{corrige}[ces ions atteignent-ils un gaz noble ?]
$n(e^-) = Z - q$ ; on compare aux nombres $2, 10, 18, 36, 54$.
\begin{enumerate}[label=\alph*)]
  \item $\ce{H^{-}}$ : $1+1 = 2$ \textbf{\textcolor{excol}{$\to$ oui}}, configuration de l'hélium.
  \item $\ce{Na^{+}}$ : $11-1 = 10$ \textbf{\textcolor{excol}{$\to$ oui}}, configuration du néon.
  \item $\ce{Cu^{2+}}$ : $29-2 = 27$ \textbf{\textcolor{attncol}{$\to$ non}} (entre $18$ et $36$) :
        le cuivre devrait perdre bien plus d'électrons pour rejoindre l'argon.
  \item $\ce{Ba^{+}}$ : $56-1 = 55$ \textbf{\textcolor{attncol}{$\to$ non}} : c'est $\ce{Ba^{2+}}$
        ($54$ électrons, configuration du xénon) qui atteindrait un gaz noble.
\end{enumerate}
Seuls $\ce{H^{-}}$ et $\ce{Na^{+}}$ conviennent.
\end{corrige}

\begin{corrige}[vrai ou faux — synthèse]
\begin{enumerate}[label=\alph*)]
  \item \textbf{\textcolor{attncol}{FAUX}} : l'hélium a $2$ électrons de valence (un \textbf{duet}),
        pas un octet.
  \item \textbf{\textcolor{excol}{VRAI}} : $2$ électrons de valence perdus $\to$ cation de charge $+2$.
  \item \textbf{\textcolor{attncol}{FAUX}} : l'électronégativité augmente vers le \emph{haut} ; c'est le
        \emph{rayon atomique} qui augmente vers le bas.
  \item \textbf{\textcolor{excol}{VRAI}} : $\ce{Cl^{-}}$ a $18$ électrons,
        $(\mathrm{K})^2(\mathrm{L})^8(\mathrm{M})^8$, comme l'argon.
\end{enumerate}
\end{corrige}

\begin{corrige}[prévoir la formule d'un composé ionique]
Chaque métal perd ses électrons de valence ; chaque non-métal complète son octet. On équilibre ensuite
les charges pour un composé neutre.
\begin{enumerate}[label=\alph*)]
  \item $\ce{Mg^{2+}}$ et $\ce{Cl^{-}}$ : il faut deux $\ce{Cl^{-}}$ pour compenser $+2$ $\to$ \ce{MgCl2}.
  \item $\ce{Na^{+}}$ et $\ce{O^{2-}}$ : il faut deux $\ce{Na^{+}}$ pour compenser $-2$ $\to$ \ce{Na2O}.
  \item $\ce{Al^{3+}}$ et $\ce{O^{2-}}$ : on équilibre $2\times(+3) = 3\times(-2)$ $\to$ \ce{Al2O3}.
\end{enumerate}
\end{corrige}

\begin{corrige}[raisonner sur les tendances]
\begin{enumerate}[label=\alph*)]
  \item \textbf{$\ce{F}$} : dans la même colonne (groupe 17), l'électronégativité augmente vers le haut,
        et $\ce{F}$ est au-dessus de $\ce{Cl}$.
  \item \textbf{$\ce{Na}$} : dans la même période, le rayon augmente vers la gauche ; $\ce{Na}$ est à
        gauche de $\ce{Mg}$.
  \item \textbf{$\ce{Li}$} : dans la même colonne (groupe 1), l'énergie d'ionisation augmente vers le
        haut, et $\ce{Li}$ est au-dessus de $\ce{Na}$.
\end{enumerate}
\end{corrige}

\end{document}
